% Shared preamble for the three worksheets (compile with xelatex).
\documentclass[12pt]{article}
\usepackage[letterpaper,left=0.5in,right=0.5in,top=0.4in,bottom=0.4in,includefoot=false,nohead]{geometry}
\usepackage{fontspec}
\setmainfont{Andika}
\usepackage{xcolor}
\usepackage{tikz}
\usetikzlibrary{arrows.meta,calc}
\usepackage{array}
\usepackage{enumitem}
\usepackage[most]{tcolorbox}
\setlist{nosep,leftmargin=1.4em}
\pagestyle{empty}
\setlength{\parindent}{0pt}
\setlength{\parskip}{0pt}
\hyphenpenalty=10000 \exhyphenpenalty=10000 \sloppy
\makeatletter\g@addto@macro\@parboxrestore{\raggedright}\makeatother
\AtBeginDocument{\raggedright}

% ---- pattern block colours ----
\definecolor{pbgreen}{HTML}{2E9E48}
\definecolor{pbblue}{HTML}{2160C4}
\definecolor{pbred}{HTML}{D7302F}
\definecolor{pbyellow}{HTML}{F2C318}
\definecolor{pbpurple}{HTML}{7B3FA0}
\definecolor{ribbontext}{HTML}{A3201C}
\definecolor{ribbontextB}{HTML}{1E6B32}
\definecolor{inkgray}{HTML}{555555}

\tikzset{
  gridline/.style={line width=0.5pt,draw=black!32},
  dotgrid/.style={line width=0.9pt,draw=black!35,dash pattern=on 1.2pt off 3pt},
  boardline/.style={line width=1.8pt,draw=black},
  holefill/.style={fill=black!78},
  upfill/.style={fill=black!20},
  recgrid/.style={line width=0.35pt,draw=black!28},
  recline/.style={line width=1pt,draw=black!85},
  tileedge/.style={line width=0.9pt,draw=black!85},
  tileoutline/.style={line width=1.5pt,draw=black},
  hlfill/.style={fill=pbyellow!60},
  ribbonfill/.style={fill=pbred!28},
  ribbonfillB/.style={fill=pbgreen!35},
  ribbon arrow/.style={-{Stealth[length=5pt]},line width=1.3pt,draw=ribbontext},
  flip arrow/.style={{Stealth[length=6pt]}-{Stealth[length=6pt]},line width=1.2pt},
  piece/.style={line width=0.8pt},
  iconline/.style={line width=0.5pt,draw=black!80},
  cubeL/.style={fill=black!62},
  cubeR/.style={fill=black!22},
  cubeS/.style={fill=white},
}

% ---- figures and icons ----
\newcommand{\fig}[1]{\input{fig/#1}}
\newcommand{\ic}[1]{\raisebox{-0.12em}{\input{fig/icon_#1}}}
\newcommand{\icb}[1]{\raisebox{-0.2em}{\input{fig/iconbig_#1}}}
\newcommand{\up}{\raisebox{0.02em}{\begin{tikzpicture}[x=1in,y=1in,line cap=round,line join=round,sharp corners]
\fill[upfill] (0,0) -- (0.2,0) -- (0.1,0.1732) -- cycle;
\draw[iconline] (0,0) -- (0.2,0) -- (0.1,0.1732) -- cycle;
\end{tikzpicture}%
}}
\newcommand{\dn}{\raisebox{0.02em}{\begin{tikzpicture}[x=1in,y=1in,line cap=round,line join=round,sharp corners]
\fill[white] (0.1,0) -- (0,0.1732) -- (0.2,0.1732) -- cycle;
\draw[iconline] (0,0.1732) -- (0.1,0) -- (0.2,0.1732) -- cycle;
\end{tikzpicture}%
}}
\newcommand{\UP}{\textbf{UP}~\up}
\newcommand{\DOWN}{\textbf{DOWN}~\dn}

% ---- answer spaces ----
\newcommand{\ansbox}[1][0.6in]{\tikz[baseline=0in]\draw[line width=1.2pt,rounded corners=3pt] (0,-0.12in) rectangle (#1,0.38in);}
\newcommand{\blank}[1][1.2in]{\rule[-0.15em]{#1}{0.6pt}}
\newcommand{\writelines}[2][7.5in]{% #2 ruled lines, 0.34in apart
  \par\noindent\begin{tikzpicture}[x=1in,y=1in]
    \useasboundingbox (0,0.02) rectangle (#1,{-0.34*#2-0.04});
    \foreach \k in {1,...,#2}{\draw[line width=0.5pt,black!70] (0,-0.34*\k) -- (#1,-0.34*\k);}
  \end{tikzpicture}\par}

% ---- page header ----
\newcommand{\bandname}{}
\newcommand{\bandcolor}{pbblue}
\newcounter{sheet}
\newcommand{\sheethead}[1]{%
  \stepcounter{sheet}%
  \noindent
  \begin{tikzpicture}[x=1in,y=1in]
    \useasboundingbox (0,-0.62) rectangle (7.5,0.24);
    \node[anchor=west,font=\small,text=inkgray] at (0,0.08) {Bellingham Math Circle \,·\, Week 1: Tiling with pattern blocks};
    \node[anchor=east,font=\small\bfseries,text=white,fill=\bandcolor,rounded corners=6pt,inner sep=4pt] at (7.5,0.08) {\bandname};
    \draw[line width=0.6pt,draw=\bandcolor] (0,-0.08) -- (7.5,-0.08);
    \node[anchor=west,circle,fill=\bandcolor,text=white,font=\large\bfseries,inner sep=2pt,minimum size=0.36in] (n) at (0,-0.38) {\thesheet};
    \node[anchor=west,font=\LARGE\bfseries] at ($(n.east)+(0.08,0)$) {#1};
  \end{tikzpicture}\par\vspace{0.14in}}
\newcommand{\nameline}{\hfill{\large Name: \blank[2.2in]}\par\vspace{0.05in}}

% actual-size note with a 1-inch check bar
\newcommand{\actualsize}[1][Boards are actual size]{%
  \par\vfill\noindent
  {\footnotesize\color{inkgray}#1: each small triangle side is 1~inch. Print at 100\%.\hfill
   This bar is 1~inch:\,\tikz[baseline=-0.5ex]{\draw[line width=0.8pt] (0,-0.05in)--(0,0.05in) (0,0)--(1in,0) (1in,-0.05in)--(1in,0.05in);}}\par}

% boxes
\newcommand{\boxit}[3][\linewidth]{%
  \par\noindent\begin{tcolorbox}[width=#1,colback=\bandcolor!6,colframe=\bandcolor,boxrule=1pt,arc=6pt,
     left=5pt,right=5pt,top=4pt,bottom=4pt,boxsep=0pt,before skip=0pt,after skip=0pt]
     \raggedright\textbf{#2} #3
  \end{tcolorbox}\par}
\newcommand{\talk}[2][\linewidth]{\boxit[#1]{Talk about it.}{#2}}
\newcommand{\lab}[1]{\tikz[baseline=-0.6ex]\node[circle,draw=\bandcolor,line width=1.2pt,text=\bandcolor,font=\bfseries,inner sep=1.5pt,minimum size=0.3in]{#1};}

\renewcommand{\bandname}{Grades 2--3}
\renewcommand{\bandcolor}{pbred}
\newcommand{\possimp}{\textbf{Possible}\quad\textbf{Impossible}}
\newcommand{\updn}{\up~\blank[0.5in]\qquad\dn~\blank[0.5in]}
\renewcommand{\arraystretch}{1.0}

\begin{document}\large

% ===================================================== 1
\sheethead{Ups and downs}
\nameline
On our boards, every small triangle points \UP{} or \DOWN.

\vspace{0.12in}
\textbf{1.} Put a dot in every \UP{} triangle on this board.

\vspace{0.15in}
\begin{center}\fig{g23_para}\end{center}

\vspace{0.1in}
\textbf{2.} Put one block at a time on the board, lined up with the lines.
How many \up{} and how many \dn{} does it cover? Try different spots and turns!

\vspace{0.12in}
\begin{center}
\setlength{\extrarowheight}{6pt}
\begin{tabular}{|>{\centering\arraybackslash}m{1.3in}|>{\centering\arraybackslash}m{1.5in}|>{\centering\arraybackslash}m{1.5in}|}
\hline
\textbf{Block} & \textbf{covers} \up & \textbf{covers} \dn \\[4pt]\hline
\icb{green} & & \\[3pt]\hline
\icb{blue} & & \\[3pt]\hline
\icb{red} & & \\[3pt]\hline
\icb{yellow} & & \\[3pt]\hline
\icb{purple} & & \\[3pt]\hline
\end{tabular}
\end{center}

\vspace{0.1in}
\talk{Did any block give two different answers? Can you place a blue so that it covers
two \up{} triangles?}
\actualsize[Board is actual size]
\newpage

% ===================================================== 2
\sheethead{Possible or impossible?}
Cover a board with \textbf{blue rhombi only} \ic{blue}: no gaps, no overlaps,
nothing hanging over the edge. Circle \textbf{Possible} or \textbf{Impossible}.
Then count the \up{} and \dn{} triangles on the board.

\vspace{0.2in}
\noindent
\begin{minipage}[t]{3.5in}
\lab{A}\hspace{0.15in}\raisebox{-1.6in}{\fig{g23_boardA}}\\[10pt]
\possimp\\[6pt]\updn

\vspace{0.22in}
\lab{B}\hspace{0.15in}\raisebox{-1.6in}{\fig{g23_boardB}}\\[10pt]
\possimp\\[6pt]\updn

\vspace{0.22in}
\lab{C}\hspace{0.15in}\raisebox{-1.6in}{\fig{g23_boardC}}\\[10pt]
\possimp\\[6pt]\updn
\end{minipage}\hfill
\begin{minipage}[t]{3.6in}
\lab{D}\hspace{0.15in}\raisebox{-3.35in}{\fig{g23_boardD}}\\[10pt]
\possimp\\[6pt]\updn

\vspace{0.4in}
\talk{Every board has an even number of triangles.
What do the impossible boards have in common?
Why does that make them impossible?}
\end{minipage}
\actualsize
\newpage

% ===================================================== 3
\sheethead{Broken hexagons}
The black triangles are \textbf{holes}: no blocks go there. Try to cover all the white
triangles with \textbf{blue rhombi only} \ic{blue}.

\vspace{0.12in}
\noindent
\begin{minipage}[c]{4.15in}\fig{g23_brokenX}\end{minipage}\hfill
\begin{minipage}[c]{3.1in}
\lab{X}\; \textbf{Hexagon X}\\[12pt]
\possimp\\[14pt]
\updn
\end{minipage}

\vspace{0.15in}
\noindent
\begin{minipage}[c]{4.15in}\fig{g23_brokenY}\end{minipage}\hfill
\begin{minipage}[c]{3.1in}
\lab{Y}\; \textbf{Hexagon Y}\\[12pt]
\possimp\\[14pt]
\updn

\vspace{0.3in}
Both hexagons have 22 white triangles. Why can one of them \emph{never} be covered?
\writelines[3.1in]{3}
\end{minipage}

\vspace{0.2in}
\boxit{Make your own.}{Which two triangles could you take out of a whole hexagon
so that it is \emph{still} possible? So that it is impossible?}
\actualsize
\newpage

% ===================================================== 4
\sheethead{Fewest gaps}
Can you cover these big triangles with blues alone? If not, put in as \textbf{many} blues
\ic{blue} as you can, then fill the gaps with greens \ic{green}.
Try to use as \textbf{few} greens as possible!

\vspace{0.25in}
\noindent
\begin{minipage}[b]{3.05in}\centering
\fig{g23_tri3}\\[0.15in]
Side 3:\quad \ic{blue} \blank[0.45in]\quad \ic{green} \blank[0.45in]
\end{minipage}\hfill
\begin{minipage}[b]{4.1in}\centering
\fig{g23_tri4}\\[0.15in]
Side 4:\quad \ic{blue} \blank[0.45in]\quad \ic{green} \blank[0.45in]
\end{minipage}

\vspace{0.3in}
\noindent
\begin{minipage}[t]{3.95in}\vspace{0pt}
\textbf{Your best so far.} (Board A on page 2 is the side-2 triangle.)

\vspace{0.12in}
\setlength{\extrarowheight}{5pt}
\begin{tabular}{|c|c|c|}
\hline
\textbf{Side} & \textbf{most blues} \ic{blue} & \textbf{fewest greens} \ic{green}\\\hline
2 & & \\\hline
3 & & \\\hline
4 & & \\\hline
\end{tabular}
\end{minipage}\hfill
\begin{minipage}[t]{3.35in}\vspace{0pt}
\talk{Compare with your group. Can anyone beat your number of greens?
Look at the triangles your greens are sitting on.
Are they \up{} or \dn?}
\end{minipage}

\vspace{0.2in}
\textbf{A plan that always works?} Describe a way to place the blues that leaves
only as many gaps as the side length. Would your plan work for \emph{every} big triangle?
\writelines{2}
\actualsize
\newpage

% ===================================================== 5
\sheethead{Bigger and bigger}
\noindent
\begin{minipage}[b]{5.05in}
\fig{g23_tri5}
\end{minipage}\hfill
\begin{minipage}[b]{2.25in}
\textbf{Side 5.} Use as many blues as you can. Fill the gaps with greens.\\[8pt]
\ic{blue} \blank[0.5in]\qquad\ic{green} \blank[0.5in]

\vspace{0.25in}
\boxit{Purple puzzle.}{Would purple pieces \ic{purple} help you leave fewer gaps?
Try covering a purple with blues.}
\end{minipage}

\vspace{0.12in}
Fill in the table. Use your boards for sides 2--5. Can you \emph{predict} sides 6 and 10?

\vspace{0.1in}
\begin{center}
\setlength{\extrarowheight}{4pt}
\begin{tabular}{|c|c|c|c|c|}
\hline
\textbf{Side} & \textbf{how many} \up & \textbf{how many} \dn & \textbf{most blues} \ic{blue} & \textbf{fewest greens} \ic{green}\\\hline
2 & & & & \\\hline
3 & & & & \\\hline
4 & & & & \\\hline
5 & & & & \\\hline
6 & & & & \\\hline
10 & & & & \\\hline
\end{tabular}
\end{center}

\vspace{0.1in}
Why can nobody ever use fewer greens than this?
(Hint: each blue covers one \up{} and one \dn.)
\writelines{2}
\actualsize[Board is actual size]
\end{document}
